\section{Design Goals}
\label{sec:navhal_design_goals}

NavHAL is designed to address the limitations of existing abstraction approaches by combining simplicity, performance, and control within a single unified framework. The design is guided by the following principles:

\begin{itemize}

    \item \textbf{Simple and Consistent API:}  
    Provide a clean and intuitive interface for peripheral interaction, enabling rapid development without exposing unnecessary hardware complexity.

    \item \textbf{Deterministic Behavior:}  
    Ensure that hardware operations have predictable and bounded execution characteristics, suitable for real-time control systems.

    \item \textbf{Low Runtime Overhead:}  
    Minimize abstraction cost by reducing indirection and leveraging compile-time configuration, ensuring performance close to direct register access.

    \item \textbf{Compile-Time Configuration:}  
    Select implementations for core, vendor, and board at build time, eliminating runtime polymorphism and reducing binary size.

    \item \textbf{Portability Across Hardware:}  
    Isolate platform-specific details while maintaining a uniform interface, enabling reuse of application code across different microcontrollers and boards.

    \item \textbf{Explicit Resource Control:}  
    Provide direct and transparent access to hardware resources without implicit allocation or hidden state management.

    \item \textbf{Execution Independence:}  
    Allow seamless operation in both bare-metal environments and with higher-level execution frameworks such as VAIOS.

    \item \textbf{Modular and Extensible Design:}  
    Enable addition of new peripherals, architectures, or boards without affecting existing implementations.

\end{itemize}

These goals collectively define a design that prioritizes predictability, efficiency, and developer usability, forming the foundation for the NavHAL architecture.